\chapter{Embolization}

\section*{Patient focus}
This chapter explains the examination or clinical problem in accessible language. Patients should use it to prepare questions and understand the purpose, experience, benefits, and limitations of imaging; local instructions and the official report take priority.


\label{ch:ch80_embolization}

Embolization is a minimally invasive procedure in which materials are deliberately introduced into a blood vessel to obstruct blood flow. It is used to control bleeding, devascularize tumours, treat vascular malformations, and reduce blood supply before surgery.

---

\section{1. Patient Perspective: What Embolization Involves}

If you have a bleeding blood vessel, a tumour that needs to be starved of its blood supply, or an abnormal tangle of blood vessels, your doctor may suggest \textbf{embolization}.

\begin{itemize}
  \item A thin tube (catheter) is inserted through a small puncture in the groin or wrist, similar to an angiogram.
  \item Under X-ray guidance, the catheter is navigated to the target blood vessel.
  \item Tiny particles, coils, or a liquid glue-like substance are injected to block the vessel.
  \item The procedure is performed under local anaesthetic with sedation. Most patients stay in hospital for 1--2 days.
  \item You may experience mild pain, fever, or nausea for a few days afterwards (post-embolization syndrome).
\end{itemize}

---

\section{2. Indications}

\begin{table}[h]
\centering
\begin{tabular}{l l}
\hline
\textbf{Indication} & \textbf{Examples} \\
\hline
Flow reduction & Preoperative embolisation of highly vascular tumours (e.g., meningioma, juvenile nasopharyngeal angiofibroma) \\
Haemorrhage control & Traumatic splenic/liver injury, postpartum haemorrhage, GI bleeding, epistaxis, haemoptysis \\
Portal vein embolisation & To induce hypertrophy of the future liver remnant before major hepatectomy \\
Tumour devascularisation & Hepatocellular carcinoma (TACE), uterine fibroids (UFE), preoperative tumour embolisation \\
Varicocele & Infertility-related varicocele embolisation \\
Vascular malformations & Arteriovenous malformations (AVMs), venous malformations, dural arteriovenous fistulae \\
\hline
\end{tabular}
\caption{Indications for Embolization}
\end{table}

---

\section{3. Embolic Agents}

The choice of embolic agent depends on the target vessel, desired endpoint, and permanence required.

\begin{table}[h]
\centering
\begin{tabular}{l l l}
\hline
\textbf{Agent} & \textbf{Type} & \textbf{Common Use} \\
\hline
Coils (detachable) & Permanent & Aneurysms, AV fistulas, vessel occlusion \\
Drug-eluting beads (DEB) & Permanent + drug & TACE for HCC (load with doxorubicin) \\
Ethanol & Permanent & Venous malformations, renal AVMs \\
Gelfoam & Temporary (weeks) & Trauma haemorrhage, preoperative embolisation \\
NBCA (glue) & Permanent & AVMs, dural AV fistulas, acute haemorrhage \\
Polyvinyl alcohol (PVA) particles & Permanent & Tumour embolisation, AVMs \\
\hline
\end{tabular}
\caption{Embolization Agents and Their Uses}
\end{table}

\begin{pearlbox}[Particle Size Matters]
The size of embolic particles determines how far they travel and what level of the vasculature they occlude. Larger particles ($>$ 500 $\mu$m) occlude proximal vessels; smaller particles ($<$ 300 $\mu$m) penetrate to the capillary level. Smaller particles carry a higher risk of non-target embolisation.
\end{pearlbox}

---

\section{4. Technique: Vessel Selection and Safety}

\begin{itemize}
  \item \textbf{Diagnostic angiography} is performed first to map the target vessel, identify collaterals, and assess for dangerous anastomoses that could lead to non-target embolisation.
  \item \textbf{Super-selective catheterisation}: The microcatheter is advanced as close as possible to the target to minimise non-target effects.
  \item \textbf{Embolisation endpoint}: Defined before starting (e.g., complete stasis, reduction in tumour blush, occlusion of a specific vessel). Over-embolisation must be avoided.
  \item \textbf{Protective measures}:
    \begin{itemize}
      \item Identify and avoid vessels supplying critical structures (e.g., spinal artery of Adamkiewicz, cystic artery, bronchial arteries supplying the spinal cord).
      \item Use coaxial microcatheter technique to maintain control.
      \item Avoid refux of embolic material by testing with contrast before injection.
    \end{itemize}
\end{itemize}

\section{5. Non-Target Embolisation}

Non-target embolisation occurs when embolic material enters vessels other than the intended target. Prevention is paramount:
\begin{itemize}
  \item Study the angiogram carefully before embolising. Identify all branches.
  \item Use the most selective catheter position possible.
  \item Inject slowly and under continuous fluoroscopic observation.
  \item If non-target embolisation occurs, immediate management may include aspiration, anticoagulation, or surgical intervention depending on the territory affected.
\end{itemize}

\begin{pearlbox}[Clinical Pearl]
In uterine fibroid embolisation (UFE), care must be taken to avoid non-target embolisation of the ovaries (via the ovarian artery) or bladder. Pre-procedure angiography should identify the uterine artery origin and any aberrant supply.
\end{pearlbox}

---

\section{6. Post-Embolization Syndrome}

Post-embolization syndrome is a common constellation of symptoms occurring 24--72 hours after the procedure:
\begin{itemize}
  \item \textbf{Fever} (low-grade, self-limiting)
  \item \textbf{Pain} at the site of embolisation (especially after splenic or hepatic embolisation)
  \item \textbf{Nausea and malaise}
  \item \textbf{Leukocytosis} (mild)
\end{itemize}
Management is supportive: analgesics, antiemetics, and IV fluids. Symptoms typically resolve within 3--5 days.

\section{7. When to Recommend Embolization vs Surgery}

\begin{itemize}
  \item \textbf{Embolisation is preferred} in haemodynamically stable trauma patients with solid organ injury (splenic, hepatic), in patients who are poor surgical candidates, in uterine fibroid management, and in highly vascular tumours where preoperative devascularisation reduces surgical blood loss.
  \item \textbf{Surgery is preferred} in haemodynamically unstable patients with active arterial bleeding unresponsive to embolisation, in contained rupture requiring definitive repair, and when embolisation is technically not feasible.
\end{itemize}

---

\section{8. Structured Reporting Example}

\begin{testbox}[Indications:]
32-year-old female, 6 weeks postpartum, with life-threatening haemorrhage unresponsive to oxytocin and uterine balloon tamponade. Haemoglobin dropping despite 4 units transfused. Angiography and embolisation requested.
\end{testbox}
\begin{refbox}[Findings:]
\textbf{Pre-embolisation angiography}: Selective catheterisation of the right uterine artery demonstrates active contrast extravasation from multiple branches supplying the lower uterine segment, consistent with active postpartum haemorrhage. The left uterine artery shows similar findings but no active extravasation.\\
\textbf{Procedure}: Using a 5 Fr Cobra catheter and 2.7 Fr microcatheter, super-selective catheterisation of the right uterine artery was performed. Gelfoam particles (1--2 mm) were injected until stasis was achieved. The left uterine artery was similarly embolised with Gelfoam. Post-embolisation angiography shows no active extravasation from either uterine artery.\\
\textbf{Complications}: None. No non-target embolisation identified.
\end{refbox}
\begin{pearlbox}[Impression:]
1. Successful bilateral uterine artery embolisation for postpartum haemorrhage.
2. Active extravasation from the right uterine artery identified and controlled with Gelfoam.
3. No complications. Haemodynamic stability expected post-procedure. Monitor haemoglobin and uterine tone.
\end{pearlbox}

\clearpage
